\documentclass[letterpaper,11pt]{article}

\usepackage[letterpaper,left=0.80in,right=0.80in,top=0.70in,bottom=0.70in]{geometry}
\usepackage[hidelinks]{hyperref}
\usepackage[english]{babel}
\input{glyphtounicode}

\pagestyle{empty}
\setlength{\parindent}{0pt}
\setlength{\parskip}{10pt}
\pdfgentounicode=1

\begin{document}

{\Large\textbf{Aryan Miriyala}}\\
\href{mailto:aryanmiriyala@gmail.com}{aryanmiriyala@gmail.com} $|$ +1 419-315-0444 $|$
\href{https://www.linkedin.com/in/aryan-miriyala-7788a21b6/}{LinkedIn} $|$
\href{https://github.com/aryanmiriyala}{GitHub}

July 17, 2026

Jane Street Software Engineering Team

Dear Jane Street Software Engineering Team,

I am applying for the Software Engineering internship because the role is built around the kind of environment I want to grow in: real production-intended projects, close mentorship from engineers, and hard programming problems where maintainability matters. I have not worked in finance or OCaml yet, but Jane Street's emphasis on learning OCaml through real team work is exactly what makes the internship interesting to me.

My strongest experience is building reliable software around messy operational systems. At AAIS, I wrote Python and JSON automation that generated more than 1,000 production SQL tables across MySQL, Oracle, PostgreSQL, and 10+ insurance lines. I also built Python/Pandas processing workflows and AWS Lambda validation that tokenized PII before moving structured output into S3. That work pushed me to care about inputs, edge cases, repeatability, and whether another engineer could trust the system later.

I have also been drawn to lower-level systems questions outside my internships. In my OpenMP runtime project, I built a C++/OpenMP experiment that recorded per-epoch telemetry and used a UCB controller to compare adaptive scheduling against fixed baselines across Mandelbrot, heat diffusion, and reduction workloads. It was the kind of project where the interesting part was not only getting code to run, but understanding why one scheduling decision behaved differently from another.

Jane Street stands out to me because the internship is honest about both sides of engineering: learning quickly and building carefully. I would bring strong Python/C++/TypeScript fundamentals, experience turning ambiguous requirements into working systems, and the humility to ask good questions while ramping on OCaml and Jane Street's internal tools.

Sincerely,\\
Aryan Miriyala

\end{document}
